\section{Hidden Markov Model：从观测反推状态}
\label{sec:06-Random-Processes/04-Applications/04}

你有一个 Markov 链在幕后运行，状态 \(Z_t\) 随着时间转移——但你根本看不见状态。你能看到的只有一串带噪声的信号 \(Y_t\)，它们依赖于当前状态，但不会直接告诉你状态是什么。

数据摆在面前：看到有人带伞出门的频率时高时低，但天气晴雨没人记录；手上有语音频谱，但音素序列未知；仪器按时给出某个读数，但生物状态隐藏在后面。你的任务是从观测序列倒推隐藏的状态序列。

这就是 Hidden Markov Model（HMM）要解决的问题。它由三条假设撑起来：

\begin{enumerate}
\tightlist
\item
  隐状态 \(Z_t\) 自己构成 Markov 链——下一步去哪，只取决于当前在哪，和更早的轨迹无关；
\item
  给定当前隐状态 \(Z_t\)，观测 \(Y_t\) 与其他所有变量条件独立——观测只''看向''当前状态，不看前后状态也不看邻居观测；
\item
  观测分布只由当前状态决定——不同状态可能产生相同的观测值（歧义），同一状态也可能产生不同的观测值（噪声）。
\end{enumerate}

\subsection{联合概率：乘积形式，但路径太多}

设初始分布 \(\pi_i=P(Z_1=i)\)，状态转移概率

\[
A_{ij}=P(Z_t=j\mid Z_{t-1}=i),
\]

以及发射概率

\[
B_j(y)=P(Y_t=y\mid Z_t=j).
\]

给定隐藏状态序列 \(z_{1:T}\) 和观测序列 \(y_{1:T}\)，联合概率由链式法则配合条件独立假设逐时间步展开：

\[
P(z_{1:T},y_{1:T})
=\pi_{z_1}B_{z_1}(y_1)
\prod_{t=2}^{T}
A_{z_{t-1}z_t}B_{z_t}(y_t).
\]

公式写出来干净，但你可不能真去枚举。隐藏状态有 \(K\) 种可能取值，\(T\) 个时刻的隐藏路径总数是 \(K^T\)。哪怕 \(K=3\)（三种天气）、\(T=100\)（一百天的观测），\(3^{100}\) 条路径——太阳年龄尺度内都跑不完。

HMM 的实用价值不在于这个联合概率公式，而在于 Markov 性质加条件独立允许你把指数级的枚举变成多项式时间的递推。以下是递推的核心。

\subsection{前向算法：观测序列出现的可能性}

问题一：给定模型参数 \(\pi,A,B\) 和观测序列 \(y_{1:T}\)，这串观测出现的概率是多少？

定义\textbf{前向量}：

\[
\alpha_t(j)=P(Y_{1:t}=y_{1:t},\,Z_t=j).
\]

它的含义：到时刻 \(t\) 为止的观测恰好是看到的那一串，并且当前（\(t\) 时刻）隐状态恰好是 \(j\) 的联合概率。不是条件概率——没有除以什么——是联合。

初始化只需要初始分布和第一个观测：

\[
\alpha_1(j)=\pi_j\,B_j(y_1).
\]

从 \(t-1\) 推进到 \(t\)：

\[
\alpha_t(j)
=B_j(y_t)\sum_i\alpha_{t-1}(i)\,A_{ij}.
\]

怎么理解？到时刻 \(t-1\) 为止的联合概率 \(\alpha_{t-1}(i)\) 已覆盖过去全部信息，现在看下一步：从状态 \(i\)（概率 \(\alpha_{t-1}(i)\)）转移到 \(j\)（概率 \(A_{ij}\)），对 \(i\) 求和等于``到 \(t\) 时刻状态变为 \(j\) 而观测覆盖到 \(t-1\) 为止''的概率。再乘一个 \(B_j(y_t)\)，把第 \(t\) 个观测也并进来，完成一个时间步的推进。

最终，把所有可能的终结状态概率加起来：

\[
P(y_{1:T})=\sum_j\alpha_T(j).
\]

复杂度是 \(O(TK^2)\)：\(T\) 个时刻，每时刻对 \(K\) 个目标状态各做一次 \(K\) 项求和。\(K^T\) 指数降到了多项式，质的差异。

前向算法本身给出的是观测序列的对数似然，可以直接用于模型选择或参数估计（比如 Baum--Welch 里的期望步）。

\subsection{过滤、预测、平滑：三个时间方向}

有了前向递推，三种推断问题的区分自然浮现：

\begin{itemize}
\item
  \textbf{过滤}：\(P(Z_t\mid Y_{1:t})\)。你只用到当前及之前的观测来推断当前状态——在线场景，实时推理。
\item
  \textbf{预测}：\(P(Z_{t+h}\mid Y_{1:t})\)。把过滤分布沿转移矩阵向前推 \(h\) 步，不需再计算观测——因为没有未来的观测可供使用。
\item
  \textbf{平滑}：\(P(Z_t\mid Y_{1:T})\)。你站在终点 \(T\) 回望时刻 \(t\)，使用未来的观测来修正对过去的判断——离线场景，知道全部观测后再重新评估。
\end{itemize}

滤波和预测可以用前向算法直接完成。平滑则需要同时使用来自过去的证据和来自未来的证据。定义\textbf{后向消息}：

\[
\beta_t(i)=P(Y_{t+1:T}=y_{t+1:T}\mid Z_t=i).
\]

从末端倒着递推：

\[
\beta_T(i)=1,\qquad
\beta_t(i)=\sum_j A_{ij}B_j(y_{t+1})\,\beta_{t+1}(j).
\]

初始值 \(\beta_T(i)=1\) 是因为 \(t=T\) 时已经没有任何未来观测需要解释了。回头走一步：从 \(i\) 转移到各 \(j\)，在 \(t+1\) 时刻产生观测 \(y_{t+1}\)，然后继续解释 \(t+2\) 往后的观测（即 \(\beta_{t+1}(j)\)）。

前向和后向的乘积一举给出平滑边缘：

\[
P(Z_t=i\mid y_{1:T})
\propto\alpha_t(i)\,\beta_t(i).
\]

原因：给定 \(Z_t=i\)，过去观测（体现于 \(\alpha_t(i)\)）和未来观测（体现于 \(\beta_t(i)\)）条件独立。两式相乘得到 \(P(Z_t=i,y_{1:T})\)，除以 \(P(y_{1:T})\) 就是平滑概率。

平滑比过滤更准确——你用了整个序列的信息重估每一个时刻——但它不能在线完成，需要等到全部观测就绪。

\subsection{Viterbi：整条路径 vs 逐点最优}

有一个常见的陷阱。你可能会想：对每个时刻独立选择最可能平滑状态

\[
\argmax_jP(Z_t=j\mid y_{1:T}),
\]

然后把它们串成一条路径。这样做得到的是逐时刻边缘最优标签，把这些标签拼起来的路径可能根本不符合转移矩阵——它可能包含概率为零的转移，甚至跳到不允许的状态序列。边缘最优和联合最优之间的裂隙，是概率推断里反复出现的问题。

Viterbi 算法寻找整条\textbf{联合}最可能路径：

\[
\argmax_{z_{1:T}}P(z_{1:T}\mid y_{1:T}).
\]

具体操作是：把前向算法里的求和换成取最大值，同时记录取得最大值的来源路径（回溯指针）。每一步存下到当前为止最优路径及其得分，最后从末端沿回溯指针复原整条最优序列。Viterbi 遵守转移约束，给出的路径在联合概率意义上最优；逐点最优路径可能连句法合法性都不具备。

两个目标适用不同场景。语音识别和词性标注要 Viterbi 路径——你需要一段自洽的状态序列。若你只关心某个特定时刻状态的概率（比如某一帧是否含某音素），平滑边缘就够，不需要去解整条路径。

\subsection{实际计算：概率下溢不是小问题}

长序列的 \(\alpha,\beta\) 是两个方向大量小于一概率的连乘。直接计算会在几十步内下溢为浮点零——不是近似不好，是彻底断掉的零。这不是理论推导的问题，是实现的必解之题。

常用对策：

\begin{itemize}
\tightlist
\item
  每个时刻对前向量归一化并记录缩放因子（归一化后所有 \(\alpha\) 值和为一），最后在对数域累加缩放因子来恢复总似然；
\item
  在对数域全部使用 log-sum-exp（保持数值稳定，代价是每次求和稍贵）；
\item
  Viterbi 全程在对数域运算——取最大值时对数同样安全。
\end{itemize}

缩放与前向后向的推导并行不是美化，是工程必需的组成部分。Baum--Welch 做参数估计时，同样需要前向-后向概率的稳定版本。

\subsection{HMM 的模型边界和替代方案}

标准 HMM 所做的假设，每一项都限制了它能建模的现象：

\begin{itemize}
\tightlist
\item
  一阶 Markov 隐状态：给定当前状态，未来独立于过去。如果系统的记忆长度长于一步，就需要高阶 HMM 或更一般的序列模型；
\item
  给定状态后观测条件独立：当前观测只依赖于当前状态，和前后观测无关。实际信号常有时域相关性（同一音素持续多帧，频谱缓慢变化），HMM 用状态序列本身的驻留来间接捕捉，并未直接建模；
\item
  状态驻留时间隐含为 Geometric：在状态 \(i\) 连续驻留 \(k\) 步的概率是 \(A_{ii}^{k-1}(1-A_{ii})\)。如果真实驻留时间长尾或有固定期限，Geometric 假设就会明显失配。Hidden semi-Markov model 放松了这一条，允许显式驻留分布；
\item
  参数不随时间变化：转移和发射矩阵推到最后一步都不变。若过程有明显的时段性、体制切换或趋势，参数恒定的假设站不住。
\end{itemize}

当这些假设不适切时，替代方向包括：

\begin{itemize}
\tightlist
\item
  状态空间模型和 Kalman filter——连续隐状态、线性高斯观测；
\item
  Particle filter——非线性非高斯连续状态空间的序列 Monte Carlo 方法；
\item
  更一般的序列模型——直接建模 \(P(Y_t\mid Y_{1:t-1})\) 而不引入隐状态（自回归、RNN、Transformer）。
\end{itemize}

HMM 的优势是可解释性：每个隐状态可以赋予一个语义标签，转移矩阵表达了状态间的因果关系，发射分布描述状态如何产生信号。代价是假设严格。它还有一个不那么显眼的麻烦——置换不可识别性：交换所有隐藏状态的名称，观测分布不变。不同初始化可能收敛到状态顺序不同的同一个解，解释标签时必须确认每个隐状态的实际含义。
